\section{Expected Contributions and Next Steps}
\label{sec:conclusion}

The accountability gap at the center of this proposal has persisted for decades without empirical scrutiny. Observers sit in boardrooms alongside directors, absorb the same strategic deliberations, and then walk out with no fiduciary obligation to protect what they heard. Whether this arrangement matters for capital markets is ultimately an empirical question, and one that the proposed research is designed to answer. The preliminary evidence suggests it does, but the more consequential finding would be that legal accountability structures shape information flow in predictable, regulable ways. If the pattern of opposite-direction responses to independent regulatory shocks survives expanded samples and more demanding specifications, the result would speak directly to how governance design channels private information into public prices.

This question sits at the boundary of several literatures that have not previously intersected. Research on interlocking directorates has studied spillovers between public firms connected through shared directors, where both firms' information is already semi-public and Regulation Fair Disclosure constrains selective disclosure. The proposed setting is fundamentally different because the observed firm is private, its board deliberations are exempt from Reg~FD, and the observer who bridges the two boardrooms owes no fiduciary duty to the private firm. Documenting a private-to-public information channel of this kind would extend the interlock literature into terrain where information asymmetry is genuine rather than attenuated by regulatory symmetry. At the same time, the focus on \textit{pre-event} price movements distinguishes this work from studies of post-event reputation contagion and decision mimicry among connected firms. The proposed regulatory shock tests would further contribute quasi-experimental evidence of a kind that the interlock literature has generally lacked, with two independent policy changes operating in opposite directions on the same mechanism.

The regulatory implications deserve emphasis. The DOJ/FTC's decision to extend interlocking directorate rules to board observers would be consistent with the preliminary finding that same-industry observer connections generate abnormal returns at connected firms before public disclosure. The NVCA's 2025 expansion of competitive harm carve-outs, which gives companies broader authority to exclude observers from sensitive discussions, would respond to a real information channel if the proposed analysis confirms the pattern. These policy developments make the proposed research timely, and the ongoing accumulation of post-2025 data creates a natural opportunity to test regulatory effectiveness in real time.

\subsection{Planned Extensions}

The preliminary findings motivate several extensions that form the core of the proposed research plan.

\begin{enumerate}
\item \textbf{Strengthening identification.} The most serious threat to a causal interpretation is that connected firms share common exposure to the same VC fund, and the VC fund itself may drive correlated returns through portfolio-wide actions unrelated to the observer channel. Three planned tests address this concern directly. Industry $\times$ event-date fixed effects would absorb common shocks hitting all firms in the same industry around the same event, isolating the observer-specific component. A control group of director-connected firms, where the same VC has a board \textit{director} rather than an observer at the private company, would test whether the effect operates through observer-specific information access or through general VC exposure. Wild cluster bootstrap inference would address the small number of clusters in the same-industry subsamples, where conventional asymptotic standard errors may be unreliable with 11 to 28 observations.

\item \textbf{Improved return measurement.} The current market-adjusted returns are a first-pass benchmark. Fama-French or DGTW characteristic-adjusted returns would account for size, value, and momentum exposures that may confound the raw market adjustment, particularly for the growth-oriented firms typical of VC portfolios. Calendar-time portfolios would serve as an alternative to event-time CARs, addressing concerns about overlapping event windows and cross-sectional dependence.

\item \textbf{Dynamic event studies for regulatory shocks.} Year-by-year event study coefficients around the NVCA 2020 and Clayton Act 2025 breaks would allow visual inspection of pre-trend dynamics and the timing of the effect onset. This approach would make the identifying assumption of parallel pre-trends directly testable rather than relying on a single structural break estimate.

\item \textbf{Sample expansion.} Incorporating events at public observed companies would approximately triple the sample and substantially improve statistical power, particularly for the bankruptcy and M\&A subsamples where connected same-industry observations are currently limited.

\item \textbf{Mechanism and channel identification.} The preliminary results establish that information reaches prices through the observer network, but the specific transmission channel remains unidentified. I plan to investigate five candidate mechanisms.
\begin{enumerate}[label=(\alph*)]
\item \textit{Direct insider trading by the observer.} If the observer personally trades on information learned at the private firm's board, the signal should appear in their Form~4 filings at the connected public company. Preliminary Form~4 analysis is suggestive but underpowered, and expanding the sample would sharpen this test.
\item \textit{VC-level portfolio rebalancing.} The observer reports back to their VC firm, which adjusts positions across portfolio companies in anticipation of the event. This channel would be detectable through 13F institutional holdings filings showing changes in the VC fund's stake before announcements at the observed firm.
\item \textit{Strategic decision-making at the connected public firm.} The observer, serving as a director at the public firm, may influence corporate decisions such as preemptive acquisitions, hiring, or capital allocation based on what they learned in the private boardroom. This would produce real effects rather than purely financial ones and could be tested using quarterly financial data around events.
\item \textit{Information diffusion through the VC ecosystem.} The observer shares information informally with colleagues, analysts, or limited partners, who then trade. This broader channel would predict stronger effects for observers at larger, more connected VC firms and for firms with greater institutional investor overlap.
\item \textit{Common investor rebalancing.} Institutional investors holding stakes in both the private and public firms, such as the VC's limited partners or crossover funds, may rebalance based on private information flow. This would be testable through 13F data on overlapping institutional holders.
\end{enumerate}
Distinguishing among these channels is a central objective of the proposed research. The null result on abnormal volume (Appendix~A4) tentatively argues against concentrated trading by a single informed party, but more granular tests are needed.

\item \textbf{Post-shock data accumulation.} As additional post-January 2025 data accumulates, the Clayton Act effect can be estimated with greater precision. The October 2025 NVCA competitive harm provisions offer a third regulatory shock to be tested once sufficient post-event data exists.

\item \textbf{Network heterogeneity.} Analysis of whether the effect varies with observer characteristics, such as the number of simultaneous board seats, the observer's role at their VC firm, or the size of the VC fund, would illuminate which types of observers are most associated with information spillover.
\end{enumerate}

\subsection{Limitations}

The current analysis has several limitations that the proposed extensions are designed to address. The observer network, while supplemented with BoardEx and Form~4 data, captures an estimated half of all observer-to-public-firm connections. The remaining connections would increase statistical power but cannot introduce spurious effects.

The most consequential identification concern is common VC exposure. Because the observer's VC fund typically holds equity in both the private and the public firm, correlated returns could reflect portfolio-level rebalancing or shared risk factors rather than information flowing through the observer channel specifically. The planned director-connected control group and industry $\times$ event-date fixed effects are designed to address this threat, but until those tests are implemented the observer-specific interpretation remains provisional.

The same-industry samples for M\&A and bankruptcy events are small, making coefficient estimates noisy despite their statistical significance. Relatedly, the network shuffle placebo test for the M\&A Buyer subsample yields $p = 0.142$, which does not reject the null at conventional levels and warrants caution in interpreting that particular result. I do not identify the specific mechanism through which information reaches prices, though the regulatory responsiveness suggests the observer's legal position is a necessary condition. Finally, the post-January 2025 sample remains short, and the Clayton Act estimates should be treated as preliminary until additional data confirms the pattern.
